\chapter{Evaluation}
\label{ch:evaluation}

\section{Evaluation Strategy}
\label{sec:eval-strategy}

\TODO{State what is and is not evaluated, up front. Technical quality against
curated golden sets: yes. A systematic audit of the deployed system against its
documentation: yes. Learning outcomes with participants: no. Being explicit
here prevents a reader inferring an outcome claim from a quality metric.}

\section{Automated Quality Evaluation}
\label{sec:automated-eval}

\thesisfigure{f18-evaluation-setup}
  {The scorecard harness: four metrics over human-curated golden sets, with a
   pass band per metric and continuous-integration gating.}
  {fig:eval-setup}

\TODO{Describe the harness from \cref{fig:eval-setup}: golden-set case types,
the pipeline abstraction that allows a deterministic fake, the pure scoring
functions, and the injected LLM judge for open-ended synthesis quality. Note
the design pattern shared with the refusal gate in \cref{sec:anti-hallucination}
--- the assertion surface is a pure function, not an HTTP round-trip, which is
what makes it usable in continuous integration.}

\section{Results}
\label{sec:results}

\thesisfigure{f19-results-scaffold}
  {Scorecard results.}
  {fig:results}

\TODO{CRITICAL --- \cref{fig:results} is currently an EMPTY SCAFFOLD. No
evaluation run exists: the results table was never applied to the production
database and no local scorecard has been recorded. Run the harness, then fill
the four metric rows and the run metadata. Do NOT write prose interpreting
results that do not exist, and do not submit with this figure as-is: either
populate it or remove both the figure and this section.}

\section{Implementation Reality versus Documented Design}
\label{sec:reality-vs-documented}

This section reports a systematic comparison of the implementation against the
project's own documentation.

\TODO{Method first, in two or three sentences: four independent readers mapped
the codebase, every architectural claim in the repository documentation was
checked against source at revision \thesisCommit{}, and each discrepancy was
recorded with a file and line reference. Method makes it a finding rather than
an anecdote.}

\TODO{Then the ledger as a longtable: twelve documented claims that do not hold,
each with the documented claim, the implemented reality, and the evidence
location. The material is already assembled in
\code{docs/thesis/ARCHITECTURE\_PRIMER.md}, section 6 --- transfer it here,
re-verifying each row against current code first, since the codebase has moved
since the primer was written.}

\TODO{Frame the discussion around \emph{why} documentation drifts in a
single-developer project under time pressure, and what would prevent it ---
generated architecture documentation, tests that assert on documented claims, or
simply deleting archived code rather than retaining it. Avoid moralising; treat
it as an engineering-process observation.}

\subsection{Built versus Deployed}
\label{sec:built-vs-deployed}

\thesisfigure{f17-flag-reality}
  {Feature-flag reality. Every flag in the project defaults to off, so the
   deployed product is a strict subset of the implemented one.}
  {fig:flag-reality}

\TODO{The reframing that makes this chapter interesting: not "what was built?"
but "what was delivered?". Cover the three inconsistent half-states where a
frontend flag is on and its backend counterpart off, producing a visible
navigation entry whose endpoint returns 404; the structural divergence between
local and deployed configuration caused by environment files excluded from the
container image; and the cascade in which a single default empties three
database tables. Close with the observation that the production default is the
least-tested configuration, since the test suite forces two flags on.}

\subsection{The Citation Gap}
\label{sec:citation-gap}

\thesisfigure{f16-citation-gap}
  {Source attribution is computed correctly on the server and consumed by no
   reachable user interface.}
  {fig:citation-gap}

\TODO{The sharpest single finding, and a good closing note for the chapter.
Citations are computed and validated server-side; the component that renders
them with clickable slide-jump exists and is imported nowhere; the two live
student surfaces either ignore the field or substitute a label derived from the
currently visible slide. So groundedness is real at the API boundary and
invisible at the user boundary. Connect this to RQ3: a user cannot perceive a
trustworthiness mechanism that is not surfaced, which has a direct bearing on
any claim about perceived reliability.}

\section{Threats to Validity}
\label{sec:threats}

\TODO{Be thorough here; it is cheap credit. At minimum: the audit was performed
by the system's own author, so confirmation bias cuts both ways; golden sets are
small and self-authored; the LLM judge shares a model family with the system
under test; no controlled comparison against static slides exists; and single-
machine deployment means nothing is known about behaviour at scale.}

\section{Unanswered Research Questions}
\label{sec:unanswered}

\TODO{Required, not optional, if RQ2 and RQ3 remain in \cref{sec:research-questions}.
Both are perception questions needing participants, and no study was conducted.
Either agree a scope reduction with your supervisor, or contribute a concrete
study design --- population, instrument, comparison condition, and the metrics
the existing event telemetry could and could not supply. A well-designed study
that was not run is a legitimate contribution; an implied answer is not.}
